\section{What Does It Mean to Model a Data Distribution?}

The previous chapters emphasized prediction, representation, transfer, and control.
In those settings the central object was often a map from an input $x$ to an output $y$ or to a useful internal representation.
Generative modeling changes the point of view.
Instead of asking only how to predict from data, we ask how to model the data distribution itself.

\medskip

\noindent
\textbf{Chapter role.}
This opening section sets the notation and taxonomy for the chapter.
It introduces the main kinds of generative models, distinguishes likelihood evaluation from sample generation, and explains why latent variables and implicit samplers lead to different mathematical and algorithmic questions.

\subsection*{Basic notation and the chapter's main question}

Let $x\in\mathcal{X}$ denote an observed data point, for example an image, a sentence, an audio segment, or a trajectory.
A dataset is written as
\[
D = \{x_1,\dots,x_n\}.
\]
A generative model with parameters $\theta$ attempts to capture the structure of the unknown data-generating distribution by introducing one or more of the following objects:
\[
p_\theta(x), \qquad p_\theta(x,z), \qquad x\sim p_\theta, \qquad x = G_\theta(\varepsilon).
\]
These symbols already suggest that generative modeling is not a single method but a family of related viewpoints.
The main question of the chapter is therefore the following:

\begin{quote}
What mathematical object is the model learning, and what operations does that object support?
\end{quote}

Some models support exact or approximate evaluation of probabilities.
Some support efficient sampling.
Some introduce latent variables that organize the data through hidden structure.
Some provide only a simulation mechanism without a tractable density.
The rest of the chapter will develop these alternatives systematically.

\subsection*{Definitions: generative model, density model, sampler, latent model, implicit model}

\begin{definition}[Generative model]
A \emph{generative model} is a parameterized family of distributions or sampling procedures intended to approximate the distribution of observed data.
Depending on the framework, this may mean learning a density $p_\theta(x)$, a joint model $p_\theta(x,z)$ with latent variables, or a sampling map that produces synthetic observations with realistic statistics.
\end{definition}

\begin{definition}[Density model]
A \emph{density model} is a generative model that assigns a probability mass or density value to each observation $x$ through an explicit expression such as
\[
p_\theta(x).
\]
When this quantity is tractable, one can evaluate likelihoods, compare models by log-likelihood, and reason directly about probability mass.
\end{definition}

\begin{definition}[Sampler]
A \emph{sampler} is a mechanism that generates synthetic data points from a model distribution.
We write
\[
x \sim p_\theta
\]
to indicate that $x$ is sampled from the model.
The ability to sample does not by itself imply that $p_\theta(x)$ can be evaluated in closed form.
\end{definition}

\begin{definition}[Latent-variable model]
A \emph{latent-variable model} introduces a hidden variable $z\in\mathcal{Z}$ and specifies a joint distribution
\[
p_\theta(x,z) = p_\theta(x\mid z)p_\theta(z).
\]
The observed-data distribution is obtained by marginalization:
\[
p_\theta(x) = \sum_z p_\theta(x,z)
\]
for discrete $z$, or
\[
p_\theta(x) = \int p_\theta(x,z)\,dz
\]
for continuous $z$.
\end{definition}

\begin{definition}[Implicit model]
An \emph{implicit model} defines a procedure for generating data, often in the form
\[
\varepsilon \sim p(\varepsilon),
\qquad
x = G_\theta(\varepsilon),
\]
without necessarily providing a tractable formula for $p_\theta(x)$.
Such models are judged primarily through sample quality, distributional matching, or downstream utility rather than direct likelihood evaluation.
\end{definition}

These definitions organize much of modern generative modeling.
Likelihood-based models, latent-variable models, adversarial generators, normalizing flows, and diffusion models all fit into this language, but they emphasize different parts of it.

\subsection*{Likelihood evaluation versus sample generation}

Two goals in generative modeling are closely related but mathematically distinct.

\paragraph{Likelihood evaluation.}
A model supports likelihood evaluation if it can compute or approximate quantities such as
\[
p_\theta(x)
\quad \text{or} \quad
\log p_\theta(x).
\]
This is useful for density estimation, uncertainty quantification, compression, and comparison by probabilistic criteria.

\paragraph{Sample generation.}
A model supports sample generation if it can draw synthetic observations
\[
x^{\mathrm{new}} \sim p_\theta.
\]
This is useful for simulation, creative generation, imputation, and scenario construction.

\medskip

\noindent
A key lesson of the chapter is that these two capabilities do not always coincide.
A model may assign a tractable density but generate poor-looking samples, or generate highly realistic samples while providing no tractable pointwise likelihood.
Thus one must always ask which operation the model is designed to support.

\begin{example}[A density that is easy to evaluate]
Suppose
\[
p_\theta(x)=\mathcal{N}(x\mid \mu,\Sigma).
\]
Then both evaluation and sampling are straightforward.
For any $x$, one can compute $p_\theta(x)$ explicitly, and one can also generate synthetic samples by drawing from the Gaussian.
This is the idealized case in which density estimation and sampling align cleanly.
\end{example}

\begin{example}[A sampler without tractable likelihood]
Suppose
\[
\varepsilon \sim \mathcal{N}(0,I),
\qquad
x = G_\theta(\varepsilon),
\]
where $G_\theta$ is a deep neural network.
Then generating $x$ may be easy, but the induced density on $x$ can be difficult or impossible to evaluate tractably.
This is the typical situation for implicit generative models.
\end{example}

\subsection*{Latent variables and hidden structure}

One of the chapter's main principles is that complicated observations may arise from simpler hidden causes.
If $z$ denotes a latent variable, then the model is written as
\[
p_\theta(x,z)=p_\theta(x\mid z)p_\theta(z).
\]
The variable $z$ may represent class identity, geometric structure, semantic content, style, trajectory intent, or other hidden factors.
The advantage of this factorization is conceptual and computational:
modeling becomes easier if the complexity of the observations can be explained through lower-dimensional or more structured latent causes.

\medskip

\noindent
Sampling from a latent-variable model typically proceeds in two stages:
\[
z \sim p_\theta(z),
\qquad
x \sim p_\theta(x\mid z).
\]
This separates the model into a prior over hidden structure and a decoder or observation model that maps hidden structure to visible data.

\medskip

\noindent
Inference, however, runs in the opposite direction.
Given an observed $x$, one would like to understand the posterior
\[
p_\theta(z\mid x)=\frac{p_\theta(x,z)}{p_\theta(x)}.
\]
This is often much harder than sampling, and it is one of the central reasons variational methods arise later in the chapter.

\subsection*{Explicit models, latent models, and implicit samplers: three worked contrasts}

This section is intended to set notation, but it should also make the distinctions concrete.
The following three examples give a compact contrast among the main model types.

\begin{example}[Explicit density model: a Gaussian on $\mathbb{R}^d$]
Let
\[
p_\theta(x)=\mathcal{N}(x\mid \mu,\Sigma),
\qquad
\theta=(\mu,\Sigma).
\]
Then:
\begin{itemize}
    \item the model directly defines a density on observations,
    \item likelihood evaluation is tractable,
    \item sampling is tractable,
    \item there is no latent variable separate from the observed $x$.
\end{itemize}
This is the cleanest form of explicit generative modeling, but it may be too restrictive for high-dimensional structured data.
\end{example}

\begin{example}[Latent-variable model: finite mixture]
Let $z\in\{1,\dots,K\}$ and define
\[
p_\theta(z=k)=\pi_k,
\qquad
p_\theta(x\mid z=k)=\mathcal{N}(x\mid \mu_k,\Sigma_k).
\]
Then the joint model is
\[
p_\theta(x,z=k)=\pi_k\,\mathcal{N}(x\mid \mu_k,\Sigma_k),
\]
and the marginal observed-data density is
\[
p_\theta(x)=\sum_{k=1}^K \pi_k\,\mathcal{N}(x\mid \mu_k,\Sigma_k).
\]
Compared with a single Gaussian, this model introduces hidden structure.
The latent variable $z$ explains which component generated the data point.
Likelihood evaluation remains tractable, but latent inference becomes part of the problem.
\end{example}

\begin{example}[Implicit sampler: neural generator]
Let
\[
\varepsilon \sim \mathcal{N}(0,I),
\qquad
x = G_\theta(\varepsilon).
\]
Then sampling is immediate: draw $\varepsilon$ and push it through $G_\theta$.
However, unless $G_\theta$ has special structure, there may be no tractable closed-form expression for the density $p_\theta(x)$.
The model is therefore suitable when realistic generation is primary, but it is less suitable when exact likelihood evaluation is required.
\end{example}

These three examples already preview the chapter's later routes:
explicit density modeling, latent-variable inference, and sample-driven implicit generation.
Normalizing flows and diffusion models will later show that some modern methods mix these viewpoints rather than falling perfectly into one category.

\subsection*{A taxonomy of generative models}

Figure~\ref{fig:gen-taxonomy} gives a compact taxonomy that will guide the rest of the chapter.

\begin{figure}[ht]
\centering
\small
\setlength{\tabcolsep}{6pt}
\renewcommand{\arraystretch}{1.25}
\begin{tabular}{|p{3.0cm}|p{3.0cm}|p{3.1cm}|p{3.2cm}|}
\hline
\textbf{Model family} & \textbf{Core object} & \textbf{Main operation} & \textbf{Typical question} \\
\hline
Explicit density model & $p_\theta(x)$ & Evaluate likelihood and sample & How much probability does the model assign to this $x$? \\
\hline
Latent-variable model & $p_\theta(x,z)$ & Marginalize or infer $z$ & Which hidden structure explains this observation? \\
\hline
Implicit model & $x=G_\theta(\varepsilon)$ & Sample generation & Can the model generate realistic observations? \\
\hline
Conditional generative model & $p_\theta(x\mid c)$ or $x=G_\theta(c,\varepsilon)$ & Controlled generation & How do we generate an $x$ consistent with condition $c$? \\
\hline
Hybrid models & combinations of the above & Trade off tractability and realism & Which operations are preserved and which are approximated? \\
\hline
\end{tabular}
\caption{A compact taxonomy of generative models. The same model may belong to more than one row; for example, latent-variable models may also be explicit density models, and some modern architectures approximate several operations at once.}
\label{fig:gen-taxonomy}
\end{figure}

\subsection*{Conditional generation and representation}

Many generative tasks are conditional rather than unconditional.
One may wish to generate an image from text, speech from phonetic content, or a completion from a prompt.
In such cases one writes
\[
p_\theta(x\mid c)
\]
for a condition $c$.
The condition may be a class label, a text description, a partial observation, or a control variable.

\medskip

\noindent
Conditional generation is especially important in modern generative AI because the system is often evaluated by whether it can produce outputs consistent with instructions or side information.
Mathematically, however, the same distinctions remain:
one may seek a conditional density, a conditional latent-variable model, or a conditional implicit sampler.

\medskip

\noindent
This also connects the chapter to representation learning.
A useful representation may serve as the condition $c$, and conversely a generative model may learn a latent code $z$ that functions as a representation.
Thus generative modeling and representation learning are not separate topics.
They are often two views of the same internal structure.

\subsection*{What this section establishes, and what it does not}

At this stage, the mathematical claims are modest but important.
The section establishes a clean vocabulary:
\begin{itemize}
    \item $p_\theta(x)$ denotes an explicit model of observed data,
    \item $p_\theta(x,z)$ denotes a joint model with latent variables,
    \item $x\sim p_\theta$ denotes sampling from the model,
    \item $x = G_\theta(\varepsilon)$ denotes an implicit sampling mechanism.
\end{itemize}
It also establishes that likelihood evaluation and sample generation are distinct operations.

\medskip

\noindent
What the section does \emph{not} yet establish is which generative route is best.
That depends on the modeling objective, the data modality, the notion of tractability one requires, and the trade-off between statistical faithfulness and computational convenience.
The remaining sections of the chapter develop these trade-offs in detail.

\subsection*{What to remember}

\begin{itemize}
    \item A generative model may define a density, a joint model with latent variables, or a sampling procedure.
    \item The notation $p_\theta(x)$, $p_\theta(x,z)$, and $x=G_\theta(\varepsilon)$ corresponds to genuinely different modeling commitments.
    \item Likelihood evaluation and sample generation are related but not identical capabilities.
    \item Latent-variable models explain visible data through hidden structure.
    \item Implicit models may generate realistic data without offering tractable pointwise likelihoods.
\end{itemize}

\subsection*{Common confusion}

\begin{itemize}
    \item \textbf{``Generative'' does not always mean ``can compute exact likelihood.''} Some generative models are defined primarily through sampling.
    \item \textbf{``Can sample'' does not imply ``understands probability well.''} High sample quality does not automatically provide calibrated density values.
    \item \textbf{Latent variables are not required in every generative model.} They are one important route, not the only route.
\end{itemize}

\subsection*{Exercises}

\begin{enumerate}
    \item Let $p_\theta(x)=\mathcal{N}(x\mid 0,I)$ on $\mathbb{R}^d$. Explain why this model is both an explicit density model and a sampler.
    \item For a finite mixture model with $K$ components, derive the formula
    \[
    p_\theta(x)=\sum_{k=1}^K \pi_k p_\theta(x\mid z=k).
    \]
    Why is this a latent-variable model?
    \item Give an example of a task in which tractable likelihood evaluation matters more than sample realism, and another in which sample realism matters more than tractable likelihood evaluation.
    \item Compare an explicit density model and an implicit sampler on the following five criteria: tractable evaluation, sample quality, hidden structure, inference, and ease of training.
\end{enumerate}

\subsection*{Notes and references}

The basic distinction between explicit density modeling, latent-variable modeling, and implicit generation runs through the modern generative-modeling literature.
Classical probability texts supply the background for density estimation and latent variables.
Mixture models and the EM tradition develop one of the oldest latent-variable routes.
Adversarial generation emphasizes implicit sampling.
Normalizing flows restore tractable density through invertible transformations.
Diffusion models show that high-quality generation can also emerge from iterative denoising constructions.
The purpose of this chapter is to place these families in one coherent mathematical picture rather than to treat them as unrelated model brands.
